\section{Results}

% Claims C01 and C02
\subsection{Cohort and classical baselines}

After technical quality control, \StudyParticipants{} participants remained
across \StudySites{} held-out sites (\cref{tab:cohort}). The connectome elastic
net had the highest equal-site balanced accuracy among the classical models
(\ConnectomeBA{}). The corresponding values were \CombinedBA{} for connectomes
plus covariates and \CovariatesBA{} for covariates alone
(\cref{tab:baselines}). These cross-site results describe computational
prediction, not clinical diagnostic performance.

\begin{table}[t]
\centering
\caption{Classical held-out-site baselines. Equal-site balanced accuracy is the primary summary.}
\label{tab:baselines}
\begin{tabular}{lccc}
\toprule
Model & Equal-site BA & Pooled BA & Pooled AUROC \\
\midrule
Covariates-only logistic regression & 0.5652 & 0.5629 & 0.5973 \\
Connectome elastic net & 0.6401 & 0.6495 & 0.6794 \\
Connectome + covariates elastic net & 0.6336 & 0.6428 & 0.6773 \\
\bottomrule
\end{tabular}
\end{table}


% Claims C03, C04, C05 and C10
\subsection{Predictive density curves}

\begin{figure}[H]
  \centering
  \includegraphics[width=0.98\textwidth]{predictive_density_curves.png}
  \caption{Equal-site mean held-out balanced accuracy across the pre-specified
  positive-edge density grid. Error bars show 1.96 standard errors across sites
  for visualization; confirmatory inference used paired site bootstrap and
  exact sign-flip tests on the locked curve estimand.}
  \label{fig:predictive-curves}
\end{figure}

All three diffusion curves declined as density increased
(\cref{fig:predictive-curves}). The normalized curve-area difference for
learned BuNN minus GCN was \BunnGcnEstimate{} (95\% site-bootstrap CI
[\BunnGcnCILow{}, \BunnGcnCIHigh{}]; exact $p=\BunnGcnP$). Because the interval
crossed zero, we detected no predictive advantage. The upper bound also fell
below the pre-specified +\PracticalMargin{} margin. This rules out a gain of
that size for the chosen estimand and pipeline, but it neither establishes
universal equivalence nor excludes smaller gains.

Against the site-matched connectome elastic net, the learned-BuNN curve
difference was \BunnElasticEstimate{} (95\% CI [\BunnElasticCILow{},
\BunnElasticCIHigh{}]; exact $p=\BunnElasticP$). Comparisons with identity
propagation and the learned-local control also failed the positive-result rule
(\cref{tab:confirmatory}).

\begin{table*}[!t]
\centering
\caption{Pre-specified predictive contrasts. Differences are BuNN minus the named comparator.}
\label{tab:confirmatory}
\footnotesize
\begin{tabularx}{\textwidth}{Xrrr}
\toprule
Contrast & Estimate & 95\% site-bootstrap CI & Exact $p$ \\
\midrule
BuNN curve - GCN curve & -0.0096 & [-0.0366, 0.0115] & 0.5237 \\
BuNN curve - trivial-bundle curve & -0.0062 & [-0.0293, 0.0108] & 0.6813 \\
BuNN nonzero mean - learned-local & -0.0167 & [-0.0347, 0.0009] & 0.0958 \\
BuNN nonzero mean - identity & -0.0146 & [-0.0320, 0.0018] & 0.1155 \\
BuNN curve - connectome elastic net & -0.0552 & [-0.0830, -0.0275] & 0.0018 \\
\bottomrule
\end{tabularx}
\end{table*}


% Claim C06
\subsection{Common-frame representation diagnostics}

\begin{figure}[H]
  \centering
  \includegraphics[width=\textwidth]{representation_density_curves.png}
  \caption{Equal-site means for common-frame layer-2 representation
  diagnostics. Zero-density points use the matched identity or learned-local
  anchor. Raw unaligned BuNN node embeddings were not compared.}
  \label{fig:representation}
\end{figure}

The matched-anchor normalized effective-rank contrast was \RankEstimate{}
(95\% CI [\RankCILow{}, \RankCIHigh{}]), the opposite of the proposed BuNN
preservation advantage. Although the raw BuNN rank curve was higher, its
learned-local anchor already included the extra map-generator capacity. The
matched contrast did not support attributing the raw difference to inter-node
bundle transport. Dispersion and cosine similarity pointed in the same
direction (\cref{fig:representation}). None of these diagnostics identifies a
biological mechanism.

% Claims C07 and C08
\subsection{Site, seed, and weighting sensitivity}

Only \LeaveOneSitePositive{} of the \StudySites{} leave-one-site estimates was
positive, and no favorable exclusion interval lay above zero (Supplementary
Figure S1). Of the five final seeds, \SeedPositive{} point estimates favored
BuNN, but \FavorableIntervals{} favorable intervals excluded zero.
\UnfavorableSeedIntervals{} unfavorable seed interval lay below zero
(Supplementary Figure S2). Participant weighting reversed the sign of the
small BuNN-GCN point estimate; the equal-site mean and median remained
negative. The pre-specified confirmatory estimand used equal-site weighting.

% Claim C09
\subsection{Observed execution cost}

Learned BuNN used \BunnParameters{} parameters, compared with \GcnParameters{}
for GCN. Its recorded fits required \BunnFitHours{} total fit-hours, while GCN
required \GcnFitHours{}. Equal-site curve balanced accuracy was lower for BuNN
(\BunnCurveBA{} versus \GcnCurveBA{}). The timings apply to this implementation
on an A10G GPU and should not be read as universal complexity estimates.

\begin{table}[H]
\centering
\caption{Observed execution cost and predictive curve summary for the three diffusion operators. Runtime is implementation- and hardware-specific.}
\label{tab:efficiency}
\small
\begin{tabular}{lrrrrr}
\toprule
Operator & Params. & Fit h & s/fit & Peak GiB & Curve BA \\
\midrule
GCN & 5,889 & 15.09 & 20.39 & 0.373 & 0.5945 \\
Trivial-bundle diffusion & 5,889 & 25.99 & 35.12 & 0.377 & 0.5911 \\
Learned BuNN & 8,529 & 31.14 & 42.08 & 0.378 & 0.5849 \\
\bottomrule
\end{tabular}
\end{table}

